% !TEX program = xelatex
\documentclass[tikz,border=4pt]{standalone}
\usepackage{fontspec}
\setmainfont{Calibri}
\setmonofont[Scale=0.88]{Consolas}
\usepackage{xcolor}
\usetikzlibrary{arrows.meta,positioning,fit,backgrounds,calc,shapes.geometric}
\definecolor{tuya}{HTML}{B45309}
\definecolor{shelly}{HTML}{0E7C86}
\definecolor{ocpp}{HTML}{15803D}
\definecolor{inny}{HTML}{1D4ED8}
\definecolor{szary}{HTML}{6B7280}
\definecolor{nasz}{HTML}{7C3AED}
\tikzset{
  box/.style={draw=#1, fill=#1!9, rounded corners=3pt, line width=0.8pt, align=center, font=\small, inner sep=4pt, minimum height=1.25cm, minimum width=2.9cm},
  grp/.style={draw=#1, dashed, rounded corners=7pt, line width=0.8pt, inner sep=7pt},
  glab/.style={font=\scriptsize\bfseries, text=#1, anchor=south west, inner sep=2pt},
  lang/.style={font=\scriptsize, text=#1, fill=white, inner sep=1.5pt, align=center},
  link/.style={{Stealth[length=2.4mm]}-{Stealth[length=2.4mm]}, line width=1.3pt, #1},
  note/.style={font=\scriptsize, text=szary, align=left},
}
\newcommand{\legenda}[2]{%
  \node[anchor=north west, font=\scriptsize, align=left] at (#1,#2) {%
  {\color{tuya}\rule[0.5ex]{0.7cm}{1.3pt}}\ język Tuya\quad
  {\color{shelly}\rule[0.5ex]{0.7cm}{1.3pt}}\ język Shelly\quad
  {\color{ocpp}\rule[0.5ex]{0.7cm}{1.3pt}}\ OCPP\quad
  {\color{szary}\rule[0.5ex]{0.7cm}{0.8pt}}\ niezależne od OCPP\quad
  {\color{nasz}\rule{0.32cm}{0.32cm}}\ tu działa tłumacz OCPP};}
\begin{document}
\begin{tikzpicture}[x=1cm,y=1cm]

\node[box=tuya] (mcu1) at (0,2.2) {\textbf{Sterownik Q11}\\mikrokontroler\\ładuje auto};
\node[box=shelly] (mod1) at (3.6,2.2) {\textbf{Moduł X1}\\system Shelly\\+ nasz skrypt};
\draw[link=tuya] (mcu1) -- node[lang=tuya, above=1pt] {Tuya} node[lang=tuya, below=1pt] {przewód} (mod1);
\begin{scope}[on background layer]
\node[grp=tuya, fit=(mcu1)(mod1)] (q1) {};
\end{scope}
\node[glab=tuya] at (q1.north west) {Q11 nr 1};

\node[box=tuya] (mcu2) at (0,0) {\textbf{Sterownik Q11}\\mikrokontroler\\ładuje auto};
\node[box=shelly] (mod2) at (3.6,0) {\textbf{Moduł X1}\\system Shelly\\+ nasz skrypt};
\draw[link=tuya] (mcu2) -- node[lang=tuya, above=1pt] {Tuya} node[lang=tuya, below=1pt] {przewód} (mod2);
\begin{scope}[on background layer]
\node[grp=tuya, fit=(mcu2)(mod2)] (q2) {};
\end{scope}
\node[glab=tuya] at (q2.north west) {Q11 nr 2};

\node[box=tuya] (mcu3) at (0,-2.2) {\textbf{Sterownik Q11}\\mikrokontroler\\ładuje auto};
\node[box=shelly] (mod3) at (3.6,-2.2) {\textbf{Moduł X1}\\system Shelly\\+ nasz skrypt};
\draw[link=tuya] (mcu3) -- node[lang=tuya, above=1pt] {Tuya} node[lang=tuya, below=1pt] {przewód} (mod3);
\begin{scope}[on background layer]
\node[grp=tuya, fit=(mcu3)(mod3)] (q3) {};
\end{scope}
\node[glab=tuya] at (q3.north west) {Q11 nr 3};

\node[box=nasz, minimum width=3.2cm, minimum height=2.2cm] (pc) at (8.9,0) {\textbf{Mały komputer}\\u klienta\\[3pt]bramka OCPP};
\foreach \i in {1,2,3} {\draw[link=shelly] (mod\i.east) -- (pc.west);}
\node[lang=shelly] at (6.4,1.45) {Shelly,\\sieć lokalna};
\begin{scope}[on background layer]
\node[grp=inny, fit=(q1)(q3)(pc)] (site) {};
\end{scope}
\node[glab=inny] at (site.north west) {Parking firmy klienta, jedna sieć lokalna};
\node[box=ocpp, minimum width=3.0cm] (csms) at (14.4,0) {\textbf{Serwer operatora}\\CSMS};
\draw[link=ocpp] (pc) -- node[lang=ocpp, above=1pt] {OCPP} node[lang=ocpp, below=1pt] {internet} (csms);
\node[note, anchor=north west] at (-1.5,-3.9) {To jest nasz test z DevKitem, tylko na stałe: laptop zastępuje mały komputer, a serwer testowy prawdziwy operator.};
\legenda{-1.5}{-4.5}

\end{tikzpicture}
\end{document}
